\section{Expansion in small reservoirs}\label{sec:residue}

We next process a graph containing a spanning expanding subgraph.
Reserving vertices and edges to connect and close paths is already
central to the cycle decomposition argument of Conlon, Fox and
Sudakov~\cite[Section~3]{CFS2014}. Our aim here is to leave residual
graphs with arbitrarily small total vertex order, so that their
recursive cost can be paid by the later folding savings. At a scale $D$, put
\[
  \sigma=D^{0.90},\qquad \tau=D^{0.74}.
\]
Throughout this section the order $r$ lies in
$\sigma\le r\le D^{1.02}$.  All assertions for sufficiently large $D$
are uniform over the graphs and orders in this range.  Their thresholds
may depend on the fixed positive constants in the statements.

\subsection{Splitting the expanding edges}

\begin{lemma}[Edge splitting]\label{lem:edge-split}
Fix $c>0$ and a positive integer $l$.  For all sufficiently large $D$,
if an $r$-vertex graph $P$ has cut expansion at least
$h_0\ge cD^{0.90}$ and $D^{0.90}\le r\le D^{1.02}$, then its edge set
can be partitioned into $l$ spanning subgraphs, each with cut expansion
at least $h_0/(2l)$.
\end{lemma}

\begin{proof}
Assign each edge independently and uniformly to one of $l$ classes.
Fix a cut whose smaller side has size $u\ge1$, and let $e$ be its
number of edges in $P$.  In any fixed class its size is binomial with
mean $e/l\ge h_0u/l$.  By \Cref{lem:chernoff}, the probability that
this class has fewer than $h_0u/(2l)$ cut edges is at most
\[
  \exp\bigl(-h_0u/(8l)\bigr).
\]
There are at most $2\binom ru\le2r^u$ choices of a cut with smaller
side of size $u$.  Thus the probability of any failure is at most
\[
  2l\sum_{u=1}^{\floor{r/2}}
       \bigl(r\exp(-h_0/(8l))\bigr)^u=o(1).
\]
Indeed, $h_0\ge cD^{0.90}$ while $\log r\le1.02\log D$, so the base
of this geometric sum tends to zero uniformly.  A successful
assignment gives the required edge partition.  Each class is taken on
the full vertex set of $P$.
\end{proof}

\subsection{Expansion after vertex sampling}

The issue is to control every sampled cut without a union bound over
all vertex subsets. Buci\'c and Montgomery establish connectivity by
short paths through a random vertex set, using successive exposures
and well-expanding subsets~\cite[Section~4]{BM2024}. In the polynomial cut-expansion
range used here, we instead fix a small family of neighborhood unions
and show that it approximates every possible failing sampled cut.

\begin{lemma}[Vertex sampling]\label{lem:vertex-sampling}
Fix $c'>0$ and $p\in(0,1)$.  Let $P$ have order
$D^{0.90}\le r\le D^{1.02}$ and cut expansion at least
$h=c'D^{0.90}$.  Select each vertex independently with probability $p$
to form $V'$.  With probability $1-o(1)$ as $D\to\infty$, all of the
following hold:
\begin{equation}\label{eq:residue-sampling-properties}
\begin{gathered}
  |V'|\le2pr,\qquad
  |N_P(v)\cap V'|\ge ph/2\quad(v\in V(P)),\\
  P[V']\text{ has cut expansion at least }\tau=D^{0.74}.
\end{gathered}
\end{equation}
The error term is uniform over the stated graphs and orders, with
$c'$ and $p$ fixed.
\end{lemma}

\begin{proof}
The minimum degree of $P$ is at least $h$.  The order bound and all
the sampled-degree bounds fail with total probability at most
\begin{equation}\label{eq:residue-size-degree-error}
  \exp(-pr/3)+r\exp(-ph/8)=o(1)
\end{equation}
by \Cref{lem:chernoff}.  We prove the expansion assertion by first
controlling a deterministic family of cuts and then approximating
any hypothetical failing sampled cut by a member of that family.

\paragraph{A deterministic witness family.}
In a failing sampled cut, the sampled-degree bound forces most vertices
on the smaller side to have many neighbors on that side. A small
selection from that side can hit their neighborhoods while exposing few
vertices across the cut; this is why we enumerate unions of a small
number of original neighborhoods.
Set
\[
  M=C_p\log r,\qquad C_p=16/p,\qquad
  s=\ceil{4rM/h}.
\]
For sufficiently large $D$, we have $M/h<1$ and $s\le r$.
Before sampling $V'$, form the family
\[
  \mathcal U=
  \left\{\ \bigcup_{v\in S}N_P(v):
                   S\subseteq V(P),\ |S|\le s\ \right\}.
\]
Its members are unions of open neighborhoods in the original graph,
and its size satisfies
\begin{equation}\label{eq:residue-witness-count}
  |\mathcal U|
  \le\sum_{t=0}^s\binom rt
  \le(s+1)r^s
  =\exp\!\bigl(O((1+rM/h)\log r)\bigr).
\end{equation}
Repeated neighborhood unions only decrease this count.

We claim that, with probability $1-o(1)$, simultaneously for every
$U\in\mathcal U$ satisfying
$\min\{|U|,r-|U|\}\ge ph/8$, one has
\begin{equation}\label{eq:residue-witness-cut}
  e_P(U\cap V',V'\setminus U)
     \ge\frac{p^2}{4}|\partial_P U|.
\end{equation}
Fix such a $U$, and put $e=|\partial_P U|$.  Expansion gives
$e\ge ph^2/8$.  Partition its cut edges into at most $2r$ matchings:
greedily color the cut edges, observing that an edge has fewer than
$2r$ other edges incident with its endpoints.  Fix one such coloring
for each cut before sampling.  The matching classes of size less than
$e/(4r)$ contain fewer than $e/2$ edges in total.  In a remaining
matching of size $m$, the indicators that both endpoints of an edge
belong to $V'$ are independent, since its edges have disjoint
endpoints.  Its retained size is therefore binomial with mean $p^2m$.
The probability of retaining fewer than $p^2m/2$ edges is at most
\[
  \exp(-p^2m/8)
  \le\exp\bigl(-p^2e/(32r)\bigr)
  \le\exp\bigl(-p^3h^2/(256r)\bigr).
\]
When none of these matching events fails, the large classes retain
at least $(p^2/2)(e/2)=p^2e/4$ edges.  Independence between different
matching classes or between different cuts is unnecessary: a union
bound bounds the probability of any failure of
\eqref{eq:residue-witness-cut} by
\begin{equation}\label{eq:residue-witness-error}
  2r|\mathcal U|\exp\bigl(-p^3h^2/(256r)\bigr)=o(1).
\end{equation}
To verify this uniformly, the scale restrictions give
\[
  \frac{h^2}{r}\ge(c')^2D^{0.78},\qquad
  \frac rh\le(c')^{-1}D^{0.12},\qquad
  M=O_p(\log D).
\]
Consequently the positive term in the exponent coming from
\eqref{eq:residue-witness-count} is
$O_{p,c'}(D^{0.12}(\log D)^2)$, while the negative term in
\eqref{eq:residue-witness-error} has magnitude at least a fixed
positive multiple of $D^{0.78}$.

\paragraph{Approximating a failing cut.}
Fix an outcome satisfying the size and degree bounds and all of
\eqref{eq:residue-witness-cut}.  Suppose, for a contradiction, that
some $A\subseteq V'$ satisfies
\begin{equation}\label{eq:residue-bad-cut}
  0<a=|A|\le |V'|/2,
  \qquad e_P(A,V'\setminus A)<\tau a.
\end{equation}
We may assume $D$ is large enough that $\tau<ph/4$.  If $a\le ph/4$,
every vertex of $A$ has more than $ph/4$ neighbors in $V'\setminus A$,
because its sampled degree is at least $ph/2$ and it has at most
$a-1$ neighbors in $A$.  This contradicts
\eqref{eq:residue-bad-cut}.  Thus
\begin{equation}\label{eq:residue-bad-cut-size}
  a>ph/4.
\end{equation}
Let
\[
  A_{\mathrm{good}}=
    \{v\in A:|N_P(v)\cap A|\ge ph/4\}.
\]
Each vertex in $A\setminus A_{\mathrm{good}}$ has more than $ph/4$
neighbors across the sampled cut.  Hence
\begin{equation}\label{eq:residue-exceptional-vertices}
  |A\setminus A_{\mathrm{good}}|
       \le\frac{4\tau a}{ph}.
\end{equation}

We now make a separate auxiliary random choice solely to prove the
existence of a suitable set $S\subseteq A$.  Include each vertex of
$A$ independently with probability $q=M/h$.  For every
$v\in A_{\mathrm{good}}$, the probability that $S$ misses all its
neighbors in $A$ is at most
\[
  (1-q)^{ph/4}\le\exp(-pM/4).
\]
Thus the probability of any required miss is at most
$r\exp(-pM/4)\le r^{-3}<1/4$.  Here the chosen value of $C_p$
is more than sufficient, since all logarithms are base two.
Also $\E|S|=aM/h$.  If $T(S)$ denotes the number of vertices of
$V'\setminus A$ with a neighbor in $S$, then a union bound for each
such vertex gives
\[
  \E T(S)
  \le\frac Mh\sum_{v\in V'\setminus A}|N_P(v)\cap A|
  =\frac Mh e_P(A,V'\setminus A)
  <\frac{\tau Ma}{h}.
\]
Markov's inequality shows that each of the events
$|S|>4aM/h$ and $T(S)>4\tau Ma/h$ has probability at most $1/4$.
The sum of these three failure probabilities is less than one.
We can therefore choose $S$ satisfying all three desired properties:
\begin{equation}\label{eq:residue-hitting-set}
\begin{gathered}
  |S|\le4aM/h,\qquad T(S)\le4\tau Ma/h,\\
  S\cap N_P(v)\ne\varnothing
      \quad(v\in A_{\mathrm{good}}).
\end{gathered}
\end{equation}
The auxiliary choice is made after $V'$ and $A$ have been fixed; the
uniform event \eqref{eq:residue-witness-cut} was already established
for every possible resulting neighborhood union.

Put $U=\bigcup_{v\in S}N_P(v)$.  Since
$|S|\le4aM/h\le4rM/h$, we have $U\in\mathcal U$.  The vertices of
$A$ absent from $U$ are all exceptional in
\eqref{eq:residue-exceptional-vertices}, and the vertices of
$U\cap V'$ outside $A$ are counted by $T(S)$.  Consequently
\begin{equation}\label{eq:residue-cut-approximation}
  \delta:=|(U\cap V')\mathbin{\triangle}A|
   \le\frac{4\tau a}{ph}+\frac{4\tau Ma}{h}
   =O_p\!\left(\frac{\tau Ma}{h}\right)=o(a).
\end{equation}
In the last step we used
$\tau M/h=O_{p,c'}(D^{-0.16}\log D)=o(1)$.
Both sides of the cut in the \emph{original} graph are large:
\[
  |U|\ge |A|-\delta=a-\delta,
  \qquad
  r-|U|\ge |V'\setminus U|
           \ge |V'\setminus A|-\delta\ge a-\delta.
\]
For sufficiently large $D$, these are at least $a/2>ph/8$.
Applying \eqref{eq:residue-witness-cut} and the expansion of $P$
therefore gives
\begin{equation}\label{eq:residue-approximating-cut-lower}
  e_P(U\cap V',V'\setminus U)\ge \frac{p^2ha}{8}.
\end{equation}
On the other hand, changing the side of one vertex changes the size
of a cut in a simple $r$-vertex graph by at most $r-1$.  Changing
the $\delta$ vertices in \eqref{eq:residue-cut-approximation} yields
\[
  e_P(U\cap V',V'\setminus U)
       \le e_P(A,V'\setminus A)+r\delta
       <\tau a+r\delta.
\]
The ratio of this upper bound to $ha$ is at most
\begin{equation}\label{eq:residue-cut-perturbation}
  \frac\tau h+
       O_p\!\left(\frac{\tau Mr}{h^2}\right)
  =O_{p,c'}(D^{-0.16})+
       O_{p,c'}(D^{-0.04}\log D)=o(1),
\end{equation}
contradicting \eqref{eq:residue-approximating-cut-lower}.
Thus no cut \eqref{eq:residue-bad-cut} exists.  Together with
\eqref{eq:residue-size-degree-error} and
\eqref{eq:residue-witness-error}, this proves the lemma.

All uses of $o(1)$ above occur after $p$ and $c'$ have been fixed.
In particular, an arbitrarily small fixed positive sampling
probability changes the cutoff but does not change any of the
strict exponent inequalities.
\end{proof}

\subsection{Closing paths through three reservoirs}

Buci\'c and Montgomery use three vertex classes and reserved connecting
subgraphs to close paths through a class avoided by the original
path~\cite[Section~2.2 and Lemma~23]{BM2024}. We use three disjoint reservoirs of
arbitrarily small fixed relative order. Every edge has both endpoints
outside at least one of them. After reserving the connecting edges,
we can therefore assign each remaining edge to such a complement,
partition that graph into paths, and close the paths through its
reservoir. The preceding sampling lemma supplies the internal expansion
needed for these routes. Unused reserved edges between a reservoir and
its complement are taken singly, leaving only unused reserved edges
inside reservoirs for recursion.

\begin{lemma}[Small residue]\label{lem:small-residue}
Fix $\eta>0$ and $c>0$.  For all sufficiently large $D$, let $G$ be
an $r$-vertex graph with
$D^{0.90}\le r\le D^{1.02}$ containing a spanning subgraph of cut
expansion at least $cD^{0.90}$.  Then $E(G)$ has a partition into at
most $9r$ completed parts, each a simple cycle or a single edge,
and the edge sets of at most three residual simple graphs.  The
residual graphs have pairwise disjoint vertex sets contained in
$V(G)$, their total order is at most $\eta r$, and each has order
strictly less than $r$.
\end{lemma}

\begin{proof}
Write $c_0=\min\{c,1\}$, and fix
\[
  p=\min\{\eta/12,1/12\},\qquad
  h=(c_0/6)\sigma.
\]
In particular, $3p<1$, $6p\le\eta$, and $2p<1$.
By \Cref{lem:edge-split}, the edges of the given spanning expanding
subgraph can be partitioned into three spanning subgraphs
$P_1,P_2,P_3$, each of cut expansion at least $h$.

\paragraph{Disjoint reservoirs and bounded representative loads.}
Label the vertices independently, using each label $i\in\{1,2,3\}$
with probability $p$ and a fourth label with probability $1-3p$.
Let $V_i$ be the vertices with label $i$.  These three sets are
pairwise disjoint, and each marginal $V_i$ is a Bernoulli-$p$
sample for the fixed graph $P_i$.  By \Cref{lem:vertex-sampling}
and a union bound over the three marginal failure probabilities,
we can choose the labels so that, for each $i$,
\begin{equation}\label{eq:residue-three-reservoirs}
\begin{gathered}
  |V_i|\le2pr,\qquad
  |N_{P_i}(v)\cap V_i|\ge ph/2\quad(v\in V(G)),\\
  P_i[V_i]\text{ has cut expansion at least }\tau.
\end{gathered}
\end{equation}
Independence between these three success events is not required.
The neighbor assertion makes $V_i$ nonempty and gives every vertex
of $V_i$ a neighbor in $V_i$, so $|V_i|\ge2$.

For each $i$, reserve all edges
$E_i^{\mathrm{in}}=E(P_i[V_i])$.  Also reserve two edges from each
$v\in V(G)\setminus V_i$ to two distinct representatives in $V_i$.
The representatives can be chosen so that at most
\begin{equation}\label{eq:residue-representative-capacity}
  k_0=\ceil{8r/(ph)}
\end{equation}
of these edges are incident with any one vertex of $V_i$.
Indeed, process the two choices for every outside vertex in turn.
At any point fewer than $2r$ edges have been chosen, so the number
of representatives already at capacity $k_0$ is at most
$2r/k_0\le ph/4$.  There are at least $ph/2$ potential
representatives for the next choice by
\eqref{eq:residue-three-reservoirs}.  Excluding the saturated ones
and, for the second choice at a vertex, its first representative,
leaves at least $ph/4-1>0$ choices for sufficiently large $D$.
This greedy procedure proves the claim.  Denote the selected
outside-to-representative edges by $E_i^{\mathrm{out}}$.

For a fixed $i$, the internal and outside reservations are disjoint.
Reservations for different values of $i$ also have disjoint edge
sets, because they belong to the edge-disjoint graphs $P_i$.
Let
\[
  E_{\mathrm{res}}=
     \bigcup_{i=1}^3(E_i^{\mathrm{in}}\cup E_i^{\mathrm{out}})
\]
be the entire reserved edge set.

\paragraph{The complementary path partitions.}
Every edge in $E(G)\setminus E_{\mathrm{res}}$ has both endpoints
outside at least one $V_i$, since its two endpoints meet at most two
of the three disjoint reservoirs.  Assign each such edge to one
such index $i$, choosing just one index if several are available.
This gives edge-disjoint graphs $G_i$ on $V(G)\setminus V_i$ whose
edge sets partition all nonreserved edges.
By \Cref{cor:endpoint-paths}, each $G_i$ has a partition into
positive-length simple paths, with each vertex an endpoint at most
twice.  Since each path has two endpoint occurrences, its path
family $\mathcal P_i$ has size at most $|V(G)\setminus V_i|\le r$.
An empty edge set contributes no paths.

For each endpoint occurrence of $v$ in $\mathcal P_i$, assign one
of its two reserved edges in $E_i^{\mathrm{out}}$, using different
edges for different occurrences.  This is possible by the endpoint
load bound.  For a path $Q\in\mathcal P_i$ with distinct endpoints
$x,y$, let $x',y'\in V_i$ be their assigned representatives.
When $x'\ne y'$, issue one demand between them in $P_i[V_i]$.
When $x'=y'$, no routing demand is necessary.
Every demand-end occurrence uses a different reserved edge at its
representative.  Hence the total demand load at each vertex of
$V_i$ is at most $k_0$.

\paragraph{Routing and simple cycles.}
Apply \Cref{lem:routing} separately in each $P_i[V_i]$, with
expansion parameter $\tau$, endpoint load $k=k_0$, singleton teams
($b=1$), and no forbidden vertices ($z=0$).  We check its numerical
hypothesis using the actual order $r_i=|V_i|$.  This graph has order
at least two and cut expansion at least $\tau$, and therefore
$\tau\le r_i-1$.  The value $\ell_i$ in \Cref{lem:routing} satisfies
\[
  \ell_i=O\!\left(\frac r\tau(\log r)^2\right),
  \qquad k_0=O_{p,c}\!\left(\frac rh\right).
\]
Here the ceiling in \eqref{eq:residue-representative-capacity}
is absorbed because $r/h\ge6$.  Consequently
\begin{equation}\label{eq:residue-routing-margin}
\begin{aligned}
  \ell_i^2(k_0+1)
   &=O_{p,c}\!\left(
        \frac{r^3}{\tau^2h}(\log r)^4\right)\\
   &=O_{p,c}(D^{0.68}(\log D)^4)
     =o(D^{0.74})=o(\tau).
\end{aligned}
\end{equation}
Thus \eqref{eq:routing-condition} holds for all sufficiently large
$D$.  The lemma supplies pairwise edge-disjoint simple routing
paths for all the distinct-representative demands in this reservoir.
All their edges are in the reserved set $E_i^{\mathrm{in}}$.

For $Q$ as above with $x'\ne y'$, join its endpoints by the edge
$xx'$, the routing path from $x'$ to $y'$, and the edge $y'y$.
Every new internal vertex lies in $V_i$, while $Q$ is entirely
outside $V_i$.  The resulting closed walk is therefore a simple
cycle.  When $x'=y'$, the two reserved edges themselves form a
simple two-edge connection from $x$ to $y$ through that
representative, and again close $Q$ to a simple cycle.  Its two
reserved edges are distinct because $x\ne y$.  In this case even a
one-edge path $Q$ produces a triangle; with different
representatives the resulting cycle has at least four edges.
Thus every completed cycle has length at least three.

The cycles are edge-disjoint.  Their original path edges partition
the nonreserved edges; each assigned representative edge is used
once; and the routing paths are edge-disjoint within each $P_i[V_i]$.
These three types of edges are disjoint by reservation.  The
reservations in different $P_i$ are also edge-disjoint, regardless
of vertex overlaps among different cycles.

\paragraph{Counting and the exact remaining edge sets.}
We have produced at most
$\sum_i|\mathcal P_i|\le3r$ cycles.  Take every unused edge in
$\bigcup_iE_i^{\mathrm{out}}$ as a single-edge part, using at most
\[
  \sum_{i=1}^3 2|V(G)\setminus V_i|\le6r
\]
further parts.  The number of completed parts is therefore at most
$9r$.

For each $i$, let $H_i$ be the graph on $V_i$ consisting of exactly
the unused edges of $E_i^{\mathrm{in}}$.  These are the only
remaining edges: every nonreserved edge was placed in a path,
every reserved outside edge was used in a cycle or taken singly,
and every reserved internal edge was used by a routing path or
placed in $H_i$.  In particular, $H_i$ is not defined by taking
all remaining edges of the original induced graph $G[V_i]$;
nonreserved edges, including ones with both endpoints in a
reservoir, were already assigned to the appropriate complementary
graphs.  This accounts for every original edge exactly once.

The graphs $H_i$ are simple, have the pairwise disjoint vertex sets
$V_i$, and satisfy
\begin{equation}\label{eq:residue-final-order}
  \sum_{i=1}^3|V(H_i)|
     =\sum_{i=1}^3|V_i|\le6pr\le\eta r,
  \qquad |V(H_i)|\le2pr\le r/6<r.
\end{equation}
Isolated vertices may be retained; a residual graph with no edges
may alternatively be omitted.  The only parameters fixed before
increasing $D$ were $\eta,c$ and the constants determined from
them.  This completes the proof.
\end{proof}
